\section{Q-Learning tabular}

\subsection{El algoritmo Q-Learning}

\begin{importantbox}{Regla de actualización de Q-Learning}
Q-Learning es un algoritmo off-policy y model-free que aprende directamente la función Q óptima:
$$Q(s, a) \leftarrow Q(s, a) + \alpha \left[r + \gamma \max_{a'} Q(s', a') - Q(s, a)\right]$$
donde:
\begin{itemize}
    \item $\alpha$: tasa de aprendizaje
    \item $r + \gamma \max_{a'} Q(s', a')$: TD target (objetivo de diferencia temporal)
    \item $r + \gamma \max_{a'} Q(s', a') - Q(s, a)$: TD error
\end{itemize}
\end{importantbox}

\begin{knowledgebox}{El carácter off-policy de Q-Learning}
Q-Learning es off-policy porque usa $\max_{a'}$ para calcular el TD target, sin importar cuál sea la política de comportamiento (la política con la que se recopilan los datos). Esto significa que podemos explorar con cualquier política (por ejemplo, $\epsilon$-greedy) y, al mismo tiempo, aprender la política óptima.
\end{knowledgebox}

\subsection{Estrategias de exploración y escalamiento de recompensas}

\begin{knowledgebox}{Estrategias de exploración comunes}
\begin{itemize}
    \item $\epsilon$-greedy: la mayor parte del tiempo elige la acción óptima y explora con probabilidad $1-\epsilon$; es adecuada para espacios de acciones discretos.
    \item Boltzmann/softmax: controla la aleatoriedad mediante una temperatura, lo que facilita tener en cuenta el valor relativo de las distintas acciones.
    \item Upper Confidence Bound (UCB): incorpora intervalos de confianza y aumenta automáticamente el peso de la exploración en los estados con pocas muestras.
\end{itemize}
\end{knowledgebox}

\begin{importantbox}{Escalamiento de recompensas y recompensas desbalanceadas}
Q-Learning es sensible a la distribución de las recompensas: una escala demasiado grande hace que el TD error se dispare, y una demasiado pequeña hace que el gradiente se desvanezca. Las prácticas habituales incluyen normalizar las recompensas, restar la reward promedio como línea base y aplicar ingeniería de forma a las recompensas dispersas (dense reward shaping). En la práctica, la instrumentation correspondiente debe registrar el TD error y la distribución de reward de cada par estado-acción, para ajustar a tiempo la estrategia de escalamiento.
\end{importantbox}

\begin{warningbox}{Convergencia engañosa con recompensas dispersas}
Cuando muchos estados nunca reciben una recompensa positiva, el agent puede converger a una política óptima local dentro de la región de recompensas negativas. Se recomienda revisar periódicamente la proporción de muestras positivas y negativas en el replay buffer, y aumentar la densidad mediante hindsight experience replay o recompensas auxiliares multitarea.
\end{warningbox}

\subsection{Condiciones de convergencia}

\begin{importantbox}{Garantía de convergencia de Q-Learning}
En el caso tabular (espacios de estados y de acciones finitos), Q-Learning tiene garantizada la convergencia a $Q^*$ si se cumplen las siguientes condiciones:
\begin{enumerate}
    \item Cada par estado--acción se visita un número infinito de veces
    \item La tasa de aprendizaje cumple las condiciones de Robbins-Monro: $\sum_t \alpha_t = \infty$, $\sum_t \alpha_t^2 < \infty$
\end{enumerate}
\end{importantbox}

\subsection{Resumen de la sección}

Q-Learning tabular aproxima directamente la función Q óptima mediante actualizaciones TD y cuenta con garantías teóricas de convergencia.
